\documentclass[11pt]{article}
\usepackage{amsmath,amssymb,amsthm}
\newtheorem{theorem}{Theorem}
\newtheorem{lemma}{Lemma}
\newtheorem{remark}{Remark}

\title{The bounded cylindric polynomials $Q_{n,c}(q)$ for $d=2$:\\
a Rogers--Ramanujan mechanism}
\author{Seed 1, Layer 2, Round 2}
\date{2026-07-04}

\begin{document}
\maketitle

\begin{abstract}
For $r=3$ and $d=2$ we prove $Q_{n,c}(q)=q^{n^2}$ for all profiles $c$ in the
cyclic orbit of $(1,1,0)$, and $Q_{n,c}(q)=q^{n(n+1)}$ for all $c$ in the orbit
of $(2,0,0)$. In particular Warnaar's positivity conjecture holds for $d=2$.
The proof runs through a general lemma, valid for every profile and rank:
multiplying the Corteel--Welsh functional equation by $(yq;q)_\infty$ converts
it into a $q$-difference system with \emph{polynomial} coefficients
$(-1)^{|J|-1}(yq;q)_{|J|-1}$. For $d=2$ this system collapses to the classical
Rogers--Ramanujan $q$-difference equation $G(y)=G(yq)+yq\,G(yq^2)$, whose
coefficients are manifestly nonnegative; this is precisely why the alternating
signs of $(yq;q)_\infty$ disappear.
\end{abstract}

\section{Setting}

Fix $r=3$. For a profile $c=(c_1,c_2,c_3)$ with $d=c_1+c_2+c_3$, let
$F_c(y,q)=\sum_{\Lambda\in\mathcal C_c}q^{|\Lambda|}y^{\max(\Lambda)}$ be the
bivariate cylindric partition generating function, and write
$F_c(y,q)=\sum_{n\ge0}f_{c,n}(q)\,y^n$. For fixed $n$ and $w$ the number of
cylindric partitions with $\max(\Lambda)=n$ and $|\Lambda|=w$ is finite (all
parts are $\le n$, so the total size $w$ bounds the number of parts); hence
each $f_{c,n}\in\mathbb Z[[q]]$ and $F_c\in\mathbb Z[[q]][[y]]$. All identities
below are identities in the ring $\mathbb Z[[q]][[y]]$, in which any element
with constant term $1$ (e.g.\ $1-yq^j$, $1-q^n$, $(yq;q)_\infty$) is a unit.

For $d=2$ we have $\ell=\gcd(d,3)=1$, so
\[
Q_{n,c}(q)=(q;q)_n\,[y^n]\bigl((yq;q)_\infty\,F_c(y,q)\bigr).
\]
Since rotating the profile rotates the cylinder --- a bijection preserving
$|\Lambda|$ and $\max(\Lambda)$ --- $F_c$, and hence $Q_{n,c}$, depends only on
the cyclic orbit of $c$. For $d=2$ there are exactly two orbits, represented
by $a=(1,1,0)$ and $b=(2,0,0)$. Write $A(y)=F_{(1,1,0)}(y,q)$ and
$B(y)=F_{(2,0,0)}(y,q)$.

Our tool is the Corteel--Welsh functional equation: for any profile $c$, with
$I_c=\{i:c_i>0\}$ and $c(J)$ the shifted profile,
\begin{equation}\label{eq:CW}
F_c(y,q)=\sum_{\emptyset\subsetneq J\subseteq I_c}(-1)^{|J|-1}\,
\frac{F_{c(J)}(yq^{|J|},q)}{1-yq^{|J|}}.
\end{equation}

\section{The key lemma: where the alternating signs go}

\begin{lemma}[$G$-transform of the Corteel--Welsh system]\label{lem:GCW}
For any rank $r$ and any profile $c$, set $G_c(y):=(yq;q)_\infty\,F_c(y,q)$.
Then
\begin{equation}\label{eq:GCW}
G_c(y)=\sum_{\emptyset\subsetneq J\subseteq I_c}(-1)^{|J|-1}\,
(yq;q)_{|J|-1}\;G_{c(J)}(yq^{|J|}).
\end{equation}
\end{lemma}

\begin{proof}
Substitute $F_{c(J)}(yq^{|J|},q)=G_{c(J)}(yq^{|J|})/(yq^{|J|+1};q)_\infty$ into
\eqref{eq:CW} and multiply both sides by $(yq;q)_\infty$. The coefficient of
$G_{c(J)}(yq^{|J|})$ becomes, with $s=|J|$,
\[
\frac{(yq;q)_\infty}{(1-yq^{s})\,(yq^{s+1};q)_\infty}
=\frac{(yq;q)_\infty}{(yq^{s};q)_\infty}
=(yq;q)_{s-1},
\]
using $(yq^{s};q)_\infty=(1-yq^{s})(yq^{s+1};q)_\infty$ and
$(yq;q)_\infty=(yq;q)_{s-1}(yq^{s};q)_\infty$.
\end{proof}

\begin{remark}
Lemma~\ref{lem:GCW} explains the role of the sign-alternating factor
$(yq;q)_\infty$ in the definition of $Q_{n,c}$: it is exactly the cofactor
that clears every denominator $1-yq^{|J|}$ from the Corteel--Welsh system,
leaving a $q$-difference system with \emph{finite polynomial} coefficients
$(yq;q)_{|J|-1}$. Since $Q_{n,c}=(q^{\ell};q^{\ell})_n\,[y^n]G_c(y)$, the
positivity conjecture is a statement about the solution of the
polynomial-coefficient system \eqref{eq:GCW}.
\end{remark}

\section{The case $d=2$}

\begin{theorem}\label{thm:main}
For all $n\ge0$,
\[
Q_{n,c}(q)=q^{n^2}\quad\text{for }c\in\{(1,1,0),(0,1,1),(1,0,1)\},
\qquad
Q_{n,c}(q)=q^{n(n+1)}\quad\text{for }c\in\{(2,0,0),(0,2,0),(0,0,2)\}.
\]
In particular the Corteel--Dousse--Uncu/Warnaar positivity conjecture holds
for $d=2$.
\end{theorem}

\begin{proof}
\emph{Step 1: the system.}
For $c=(1,1,0)$ we have $I_c=\{1,2\}$ and (indices cyclic, $c_0=c_3$):
$c(\{1\})=(0,2,0)$, $c(\{2\})=(1,0,1)$, $c(\{1,2\})=(0,1,1)$. Using cyclic
invariance to replace $F_{(0,2,0)}$ by $B$ and $F_{(1,0,1)},F_{(0,1,1)}$ by
$A$, \eqref{eq:CW} reads
\begin{equation}\label{eq:CWa}
A(y)=\frac{B(yq)}{1-yq}+\frac{A(yq)}{1-yq}-\frac{A(yq^2)}{1-yq^2}.
\end{equation}
For $c=(2,0,0)$ we have $I_c=\{1\}$ and $c(\{1\})=(1,1,0)$, so
\begin{equation}\label{eq:CWb}
B(y)=\frac{A(yq)}{1-yq}.
\end{equation}

\emph{Step 2: eliminate $B$.}
Substituting \eqref{eq:CWb} (at $y\mapsto yq$) into \eqref{eq:CWa},
\[
A(y)=\frac{A(yq)}{1-yq}
+A(yq^2)\left[\frac{1}{(1-yq)(1-yq^2)}-\frac{1}{1-yq^2}\right]
=\frac{A(yq)}{1-yq}+\frac{yq\,A(yq^2)}{(1-yq)(1-yq^2)},
\]
since $1-(1-yq)=yq$. Multiplying by $1-yq$:
\begin{equation}\label{eq:star}
(1-yq)\,A(y)=A(yq)+\frac{yq}{1-yq^2}\,A(yq^2).
\end{equation}

\emph{Step 3: the $G$-transform.}
Let $G(y)=(yq;q)_\infty A(y)$. Substituting
$A(y)=G(y)/(yq;q)_\infty$, $A(yq)=G(yq)/(yq^2;q)_\infty$,
$A(yq^2)=G(yq^2)/(yq^3;q)_\infty$ into \eqref{eq:star} and multiplying by
$(yq;q)_\infty=(1-yq)(yq^2;q)_\infty=(1-yq)(1-yq^2)(yq^3;q)_\infty$ gives
\[
(1-yq)\,G(y)=(1-yq)\,G(yq)+yq\,(1-yq)\,G(yq^2).
\]
Dividing by the unit $1-yq$:
\begin{equation}\label{eq:RR}
G(y)=G(yq)+yq\,G(yq^2).
\end{equation}
This is the classical Rogers--Ramanujan $q$-difference equation.
(Equivalently: apply Lemma~\ref{lem:GCW} to both profiles --- the $|J|=1$
terms carry coefficient $+1$, the single $|J|=2$ term carries $-(1-yq)$, and
elimination of $G_B$ recombines the signs into the nonnegative coefficients
$1$ and $yq$ of \eqref{eq:RR}.)

\emph{Step 4: uniqueness and solution.}
Write $G(y)=\sum_{n\ge0}g_n(q)y^n$. Comparing coefficients of $y^n$ in
\eqref{eq:RR}:
\[
g_n=q^n g_n+q^{2(n-1)}\cdot q\,g_{n-1}
\qquad\Longleftrightarrow\qquad
(1-q^n)\,g_n=q^{2n-1}g_{n-1}\quad(n\ge1).
\]
Since $1-q^n$ is a unit of $\mathbb Z[[q]]$, the $g_n$ are uniquely determined
by $g_0$. Now $g_0=G(0)=F_c(0,q)=1$ (only the empty cylindric partition has
maximum $0$), and by induction, using $\sum_{j=1}^{n}(2j-1)=n^2$,
\[
g_n=\frac{q^{n^2}}{(q;q)_n}.
\]
Therefore
\[
Q_{n,(1,1,0)}(q)=(q;q)_n\,g_n=q^{n^2}.
\]

\emph{Step 5: the orbit of $(2,0,0)$.}
Let $G_B(y)=(yq;q)_\infty B(y)$. By \eqref{eq:CWb},
\[
G_B(y)=\frac{(yq;q)_\infty}{1-yq}\,A(yq)=(yq^2;q)_\infty A(yq)=G(yq),
\]
so $[y^n]G_B=q^n g_n=q^{n(n+1)}/(q;q)_n$ and
$Q_{n,(2,0,0)}(q)=q^{n(n+1)}$.

By cyclic invariance the same formulas hold across each orbit. Both are
monomials, hence have nonnegative coefficients; and
$Q_{n,c}(1)=1=\bigl(\tfrac16(d+1)(d+2)-1\bigr)^n$ for $d=2$, consistent with
Welsh's evaluation.
\end{proof}

\begin{remark}[Example]
$n=2$, $c=(1,1,0)$: the recursion gives
$g_1=q/(1-q)$, $g_2=q^{3}g_1/(1-q^2)=q^{4}/\bigl((1-q)(1-q^2)\bigr)$, so
$Q_{2}=(q;q)_2\,g_2=q^{4}$. Direct expansion of
$(q;q)_2\,[y^2]\bigl((yq;q)_\infty F_{(1,1,0)}(y,q)\bigr)$, with
$f_{c,0}=1$, $f_{c,1}=2q/(1-q)$,
$f_{c,2}=q^2(2+q+q^2)/\bigl((1-q)^2(1+q)\bigr)
=2q^2+3q^3+6q^4+7q^5+10q^6+\cdots$
(machine-verified against the Corteel--Welsh system), reproduces $q^4$ after
the massive cancellation --- verified to precision $q^{900}$.
\end{remark}

\begin{remark}[Why the signs collapse, and what is special about $d=2$]
Three separate phenomena combine:
(i)~\emph{globally}, Lemma~\ref{lem:GCW} absorbs the infinite alternating
series $(yq;q)_\infty$ into finite coefficients $(yq;q)_{|J|-1}$;
(ii)~\emph{for $d=2$}, elimination of the second orbit recombines the residual
signs $\{+1,+1,-(1-yq)\}$ into the manifestly nonnegative coefficients
$\{1,yq\}$ of \eqref{eq:RR}; (iii)~the resulting recursion
$(1-q^n)g_n=q^{2n-1}g_{n-1}$ is first order with monomial coefficient, which
is why $Q_n$ degenerates to a single monomial rather than a positive
polynomial. For $d>2$, step~(iii) genuinely fails (the eliminated system has
higher order), but steps (i)--(ii) formulate the conjecture as a positivity
property of the polynomial-coefficient system \eqref{eq:GCW}; finding the
general-$d$ analogue of the sign recombination (ii) is the open problem.
\end{remark}

\begin{remark}[Relation to known results]
$t=k+d=5$ is the Rogers--Ramanujan modulus; equation \eqref{eq:RR} and its
solution are classical (Rogers 1894), and Corteel--Welsh (2019) used
\eqref{eq:CW} at $d=2$ to reprove the Rogers--Ramanujan identities. Warnaar
(2023) proved positivity for $d=2$ via explicit multisums. The content here is
the clean statement and proof in the $Q_{n,c}$ normalization, together with
Lemma~\ref{lem:GCW}, which is profile- and rank-independent and isolates
exactly where the alternating signs go.
\end{remark}

\section*{Verification}
Script \texttt{scratch/scripts/seed1\_R2L2\_d2\_verify.sage}: (1) brute-force
enumeration of all cylindric partitions of size $\le25$ for all six profiles
matches the Corteel--Welsh solution (validating \eqref{eq:CWa},
\eqref{eq:CWb} and the shifted-profile computations); (2) $Q_{n}=q^{n^2}$ and
$Q_n=q^{n(n+1)}$ verified for $n\le10$ at precision $q^{900}$ (well above the
truncation-safety bound $6n^2+200=800$); (3) equations \eqref{eq:star},
\eqref{eq:RR}, \eqref{eq:GCW}, $G_B(y)=G(yq)$ and $g_n=q^{n^2}/(q;q)_n$ each
verified coefficientwise. All checks passed.

\end{document}
